\documentclass[reqno,11pt]{amsart}
\usepackage{math327-lecture}
\usepackage{needspace}
\lectureinfo{8}{Sep. 21, 2026}{Automorphisms, Normal Subgroups, and Sign}

% Based on the mathematical content of handwritten Lecture 6, Sep. 15, 2025.
\begin{document}
\makelecturetitle

\section{Homomorphisms, isomorphisms, and automorphisms}

We begin by recalling the definitions from Lecture 7. For groups $G$ and $H$,
a map $\varphi:G\to H$ is a homomorphism if
\[
  \varphi(ab)=\varphi(a)\varphi(b)\qquad(a,b\in G).
\]
It follows that $\varphi(e_G)=e_H$ and
$\varphi(a^{-1})=\varphi(a)^{-1}$. We write
\[
  \Hom(G,H)=\{\varphi:G\to H:\varphi\text{ is a homomorphism}\}.
\]

\begin{define}[Isomorphism and automorphism]
An \emph{isomorphism} $\varphi:G\to H$ is a homomorphism admitting a
homomorphic inverse $\psi:H\to G$, so that
\[
  \psi\circ\varphi=\id_G,\qquad \varphi\circ\psi=\id_H.
\]
We write $G\cong H$ if such a map exists, and $\Iso(G,H)$ for the set of
isomorphisms from $G$ to $H$.
An isomorphism $G\to G$ is called an \emph{automorphism} of $G$; the set
of all automorphisms is denoted by $\Aut(G)$.
\end{define}

As proved in Lecture 7, for a homomorphism $\varphi:G\to H$,
\begin{align*}
  \varphi\text{ is an isomorphism}
  &\quad\Longleftrightarrow\quad \varphi\text{ is bijective}\\
  &\quad\Longleftrightarrow\quad
  \ker(\varphi)=\{e_G\}\text{ and }\im(\varphi)=H.
\end{align*}
The kernel is a special kind of subgroup, called a normal subgroup.
To explain this, we first introduce an important automorphism.

\section{Example: conjugation}

\begin{define}[Conjugation]
Fix $g\in G$. The map
\begin{equation}
  c_g:G\longrightarrow G,\qquad c_g(a)=gag^{-1},
\end{equation}
is called \emph{conjugation by $g$}.
\end{define}

If $G$ is abelian, then $c_g(a)=a$ for every $a,g\in G$, so each $c_g$
is the identity map. Conjugation becomes useful in studying nonabelian groups.

\begin{proposition}
For every $g\in G$, the map $c_g$ is an automorphism, with inverse $c_{g^{-1}}$.
\end{proposition}
\begin{proof}
For $a,b\in G$, inserting $g^{-1}g=e_G$ gives
\[
  c_g(ab)=gabg^{-1}=ga(g^{-1}g)bg^{-1}
  =(gag^{-1})(gbg^{-1})=c_g(a)c_g(b).
\]
Thus $c_g$ is a homomorphism. Moreover,
\[
  c_g\bigl(c_{g^{-1}}(a)\bigr)=g(g^{-1}ag)g^{-1}=a,
  \qquad
  c_{g^{-1}}\bigl(c_g(a)\bigr)=g^{-1}(gag^{-1})g=a.
\]
Hence $c_{g^{-1}}$ is its inverse.
\end{proof}

\begin{remark}
Conjugation distributes over any finite product:
\[
  g(a_1\cdots a_r)g^{-1}=(ga_1g^{-1})\cdots(ga_rg^{-1}).
\]
Isomorphisms, including automorphisms, preserve properties expressed in
terms of the group operation. For example, conjugate elements have the
same order, since $(gag^{-1})^m=ga^mg^{-1}$ for every integer $m$.
\end{remark}


\subsection*{Conjugation in cycle notation}

Conjugation of permutations has a particularly simple description:
it relabels the entries of each cycle.

\begin{proposition}
	For $\sigma\in S_n$ and a cycle $\tau=(j_1\ \cdots\ j_r)$,
	\begin{equation}
		\sigma\tau\sigma^{-1}
		=\bigl(\sigma(j_1)\ \cdots\ \sigma(j_r)\bigr).
	\end{equation}
\end{proposition}
\begin{proof}
	Applying the three permutations from right to left sends
	\[
	\sigma(j_i)\xmapsto{\sigma^{-1}}j_i
	\xmapsto{\tau}j_{i+1}
	\xmapsto{\sigma}\sigma(j_{i+1}),
	\]
	where $j_{r+1}=j_1$. Every entry outside
	$\{\sigma(j_1),\ldots,\sigma(j_r)\}$ is fixed. This is exactly the cycle
	on the right-hand side.
\end{proof}

Since conjugation distributes over products, the same rule applies to
every cycle in a disjoint-cycle decomposition. In particular, conjugation
preserves cycle type.


\section{Normal subgroups}

\begin{define}[Normal subgroup]
A subgroup $N\leq G$ is \emph{normal} if
\begin{equation}
  gNg^{-1}\subseteq N\qquad\text{for every }g\in G,
\end{equation}
where $gNg^{-1}=\{gng^{-1}:n\in N\}$. We write $N\trianglelefteq G$.
\end{define}

\begin{remark}
The inclusion is automatically an equality. Applying it to $g^{-1}$ gives
$g^{-1}Ng\subseteq N$. Conjugating this inclusion by $g$ gives
$N\subseteq gNg^{-1}$. Therefore $gNg^{-1}=N$ for every $g\in G$.
\end{remark}

Normal subgroups arise naturally as kernels. They will also allow us to
form quotient groups and study a group through smaller groups, using the
isomorphism theorems.

\begin{proposition}\label{prop:kernel-normal}
If $\varphi:G\to H$ is a homomorphism, then $\ker(\varphi)\trianglelefteq G$.
\end{proposition}
\begin{proof}
We already know that $\ker(\varphi)$ is a subgroup. If $a\in\ker(\varphi)$
and $g\in G$, then
\[
  \varphi(gag^{-1})
  =\varphi(g)\varphi(a)\varphi(g)^{-1}
  =\varphi(g)e_H\varphi(g)^{-1}=e_H.
\]
Thus $gag^{-1}\in\ker(\varphi)$, as required.
\end{proof}

\begin{example}
Every subgroup of an abelian group is normal, since $gag^{-1}=a$.
For any group $G$, the subgroups $\{e_G\}$ and $G$ are normal.
A normal subgroup different from both is called a \emph{proper nontrivial
normal subgroup}.
\end{example}

\begin{define}[Simple group]
A nontrivial group $G$ is \emph{simple} if its only normal subgroups are
$\{e_G\}$ and $G$.
\end{define}

\subsection*{Two cyclic subgroups of $S_3$}

Write
\[
  S_3=\{1,x,x^2,y,xy,x^2y\},\qquad
  x^3=1,\quad y^2=1,\quad yx=x^2y.
\]
Consider the cyclic subgroups
\[
  H=\langle x\rangle=\{1,x,x^2\},\qquad L=\langle y\rangle=\{1,y\}.
\]
Both are abelian as groups in their own right. Are they normal in $S_3$?

\begin{example}[The subgroup $H$ is normal]
Conjugation by any element of $H$ fixes every element of $H$.
To conjugate by $y$, use $y^{-1}=y$ and the relation $yx=x^2y$:
\begin{align*}
  yxy^{-1}&=yxy=x^2,\\
  yx^2y^{-1}&=(yxy^{-1})(yxy^{-1})=x^2x^2=x.
\end{align*}
Consequently,
\begin{align*}
  yHy^{-1}
  &=\{y,yx,yx^2\}y^{-1}\\
  &=\{1,yxy^{-1},yx^2y^{-1}\}
   =\{1,x^2,x\}=H.
\end{align*}
This also deals with the remaining elements $x^iy$, because
\[
  (x^iy)H(x^iy)^{-1}=x^i(yHy^{-1})x^{-i}=x^iHx^{-i}=H.
\]
Thus $H\trianglelefteq S_3$.
\end{example}

\begin{example}[The subgroup $L$ is not normal]
It is enough to find one conjugation that does not preserve $L$.
Since $x^{-1}=x^2$ and $xyx=y$, we have
\[
  xyx^{-1}=xyx^2=(xyx)x=yx=x^2y.
\]
Therefore
\[
  xLx^{-1}=\{1,xyx^{-1}\}=\{1,x^2y\}\neq\{1,y\}=L.
\]
In particular, a subgroup can be abelian without being normal in its
ambient group.
\end{example}

With $x=(1\ 2\ 3)$ and $y=(1\ 2)$, the six elements are
\[
\begin{array}{c|cccccc}
  g&1&x&x^2&y&xy&x^2y\\ \hline
  \text{cycle notation}&(1)&(1\ 2\ 3)&(1\ 3\ 2)&(1\ 2)&(1\ 3)&(2\ 3).
\end{array}
\]
The subgroup $H$ consists of the even permutations. We next construct the
sign homomorphism and identify $H$ as its kernel, giving another proof
that $H$ is normal.




\section{Examples: determinant and sign}

\begin{example}[The determinant homomorphism]
Multiplicativity of the determinant says that
\begin{equation}
  \det:\GL_n(\R)\longrightarrow\R^*,
  \qquad A\longmapsto\det(A),
\end{equation}
is a homomorphism, where $\R^*=\R\setminus\{0\}$ is a group under
multiplication. Its kernel is
\[
  \ker(\det)=\SL_n(\R)
  =\{A\in\GL_n(\R):\det(A)=1\}.
\]
Since kernels are normal, $\SL_n(\R)\trianglelefteq\GL_n(\R)$.
\end{example}

\subsection*{Approach A: parity of permutation length}

\begin{example}[The sign homomorphism]
Let $\{\pm1\}$ be the group under usual multiplication. Up to isomorphism,
it is the unique group with two elements.
Recall that $\ell(\sigma)$ is the minimum number of simple transpositions
needed to express $\sigma$; it also equals the number of inversions of
$\sigma$. Call $\sigma$ \emph{even} or \emph{odd} according to the parity
of $\ell(\sigma)$.

We will show that the following map is a homomorphism:
\begin{equation}
  \sgn:S_n\longrightarrow\{\pm1\},\qquad
  \sgn(\sigma)=(-1)^{\ell(\sigma)}=
  \begin{cases}
    +1,&\text{if $\sigma$ is even},\\
    -1,&\text{if $\sigma$ is odd}
  \end{cases}
  .
\end{equation}
The multiplication rule in $\{\pm1\}$ has exactly the desired parity table:
\[
\begin{array}{c|cc}
 &\text{even}&\text{odd}\\ \hline
\text{even}&\text{even}&\text{odd}\\
\text{odd}&\text{odd}&\text{even}
\end{array}
\]
\end{example}

\begin{proof}
The inversion lemma from Lecture 5 says that right multiplication by a
simple transposition $s_i=(i\ i+1)$ changes the number of inversions by
exactly one. Since length equals the number of inversions,
\[
  \ell(ps_i)=\ell(p)\pm1,
  \qquad\text{so}\qquad
  \ell(ps_i)\equiv\ell(p)+1\pmod2.
\]
Let $\sigma,\tau\in S_n$ and choose a shortest expression
$\tau=s_{i_1}\cdots s_{i_k}$, where $k=\ell(\tau)$.
Starting with $\sigma$ and multiplying successively on the right by
$s_{i_1},\ldots,s_{i_k}$ flips the parity at each step. Thus
\[
  \ell(\sigma\tau)\equiv\ell(\sigma)+k
  =\ell(\sigma)+\ell(\tau)\pmod2.
\]
It follows that
\[
  \sgn(\sigma\tau)=(-1)^{\ell(\sigma\tau)}
  =(-1)^{\ell(\sigma)+\ell(\tau)}
  =\sgn(\sigma)\sgn(\tau).
\]
\end{proof}

\subsection*{Approach B: permutation matrices}

We next construct the sign homomorphism geometrically. For
$\sigma\in S_n$, let $e_1,\dots,e_n$ be the standard basis of $\R^n$ and
define the permutation matrix $[\sigma]$ by
\begin{equation}
  [\sigma]e_i=e_{\sigma(i)}.
\end{equation}
Equivalently, the $i$th column of $[\sigma]$ is $e_{\sigma(i)}$.
If $E_{ij}$ denotes the matrix with a $1$ in row $i$, column $j$, and zeros
elsewhere, then
\begin{equation}
  [\sigma]=\sum_{i=1}^n E_{\sigma(i),i}.
\end{equation}

\begin{example}
For $\sigma=(1\ 2\ 3)\in S_3$,
\[
  [\sigma]=
  \begin{pmatrix}
	    0&0&1\\
	    1&0&0\\
	    0&1&0
	  \end{pmatrix},
  \qquad
  [\sigma]
  \begin{pmatrix}x_1\\x_2\\x_3\end{pmatrix}
  =
  \begin{pmatrix}x_3\\x_1\\x_2\end{pmatrix}.
\]

\end{example}

\begin{warningbox}
The entry in output position $j$ is $x_{\sigma^{-1}(j)}$.
Indeed, the $i$th basis vector is sent to $e_{\sigma(i)}$, so the entry
originally in position $i$ appears in position $\sigma(i)$.
In column-vector notation,
\[
  [\sigma]
  \begin{pmatrix}x_1\\x_2\\\vdots\\x_n\end{pmatrix}
  =
  \begin{pmatrix}
    x_{\sigma^{-1}(1)}\\x_{\sigma^{-1}(2)}\\\vdots\\x_{\sigma^{-1}(n)}
  \end{pmatrix}.
\]
\end{warningbox}

\begin{proposition}
The assignment $\sigma\mapsto[\sigma]$ is a homomorphism
$S_n\to\GL_n(\R)$.
\end{proposition}

\begin{proof}
Using right-to-left composition, for every standard basis vector $e_i$ we have
\[
  [\sigma][\tau]e_i
  =[\sigma]e_{\tau(i)}
  =e_{\sigma(\tau(i))}
  =[\sigma\tau]e_i.
\]
Since the matrices agree on every basis vector,
$[\sigma][\tau]=[\sigma\tau]$. Also, $[\id]=I_n$, so
\[
  [\sigma][\sigma^{-1}]=I_n=[\sigma^{-1}][\sigma].
\]
Thus $[\sigma]^{-1}=[\sigma^{-1}]$, which shows that $[\sigma]$ belongs
to $\GL_n(\R)$ and completes the proof.
\end{proof}

Both $[\sigma]$ and $[\sigma^{-1}]$ have integer entries, so their
determinants are integers. Since they are inverse matrices,
\[
  \det([\sigma])\det([\sigma^{-1}])=1.
\]
Hence $\det([\sigma])\in\{+1,-1\}$.

\Needspace{9\baselineskip}
\begin{define}[Sign]
The determinant construction of the \emph{sign homomorphism} is the composition
of $[\cdot]:S_n\longrightarrow\GL_n(\R)$, $\sigma\mapsto[\sigma]$,
with $\det:\GL_n(\R)\longrightarrow\R^*$:
\begin{equation}
  \sgn:=\det\circ[\cdot]:S_n\longrightarrow\R^*,
  \qquad \sgn(\sigma)=\det([\sigma]).
\end{equation}
Its image lies in $\{\pm1\}$, so we may regard it as a map
$S_n\to\{\pm1\}$. We verify below that it agrees with Approach A.
\end{define}

The map is a homomorphism because
\[
  \sgn(\sigma\tau)=\det([\sigma\tau])
  =\det([\sigma][\tau])=\sgn(\sigma)\sgn(\tau).
\]

This calculation is an instance of the following general fact, which is left as an exercise.

\begin{lemma}[Composition of homomorphisms]
Let $G_1,G_2,G_3$ be groups, and let $\varphi:G_1\to G_2$ and
$\psi:G_2\to G_3$ be homomorphisms. Then
$\psi\circ\varphi:G_1\to G_3$ is a homomorphism.
\end{lemma}



Finally, to compare the two constructions of sign, note that the matrix of a simple
transposition exchanges two columns of $I_n$ and has determinant $-1$.
For a shortest expression $\sigma=s_{i_1}\cdots s_{i_k}$, where
$k=\ell(\sigma)$, multiplicativity gives
\[
  \det([\sigma])=\prod_{j=1}^k\det([s_{i_j}])
  =(-1)^k=(-1)^{\ell(\sigma)}.
\]
Thus the determinant construction agrees with the sign homomorphism from
Approach A.

\subsection*{The kernel of sign}

\begin{define}[Alternating group]
The \emph{alternating group} is the subgroup of even permutations:
\[
  A_n=\ker(\sgn)=\{\sigma\in S_n:\sgn(\sigma)=1\}.
\]
By Proposition~\ref{prop:kernel-normal}, $A_n\trianglelefteq S_n$.
\end{define}

For $S_3$, we obtain
\[
  \ker(\sgn)=A_3=\{(1),(1\ 2\ 3),(1\ 3\ 2)\}
  =\{1,x,x^2\}=H.
\]
The other three elements, $(1\ 2),(1\ 3),(2\ 3)$, have sign $-1$.
This recovers the normality of $H$ from the general fact that kernels are normal.


\end{document}
